\documentclass[12pt,a4paper]{article}
\usepackage[utf8]{inputenc}
\usepackage[T1]{fontenc}
\usepackage[italian]{babel}
\usepackage{geometry}
\usepackage{hyperref}
\usepackage{booktabs}
\usepackage{parskip}
\usepackage{titlesec}
\usepackage{xcolor}
\usepackage{listings}

\geometry{margin=2.5cm}

\hypersetup{
    colorlinks=true,
    linkcolor=blue!60!black,
    urlcolor=blue!60!black,
}

\titleformat{\section}{\large\bfseries}{}{0em}{}[\titlerule]
\titleformat{\subsection}{\normalsize\bfseries\itshape}{}{0em}{}

\title{\textbf{Relazione Tecnica}\\[0.4em]
       \large Progetto: Agenti RL e LLM/VLM su MiniGrid-DoorKey}
\author{}
\date{\today}

\begin{document}

\maketitle
\tableofcontents
\bigskip

%─────────────────────────────────────────────────────────────
\section{Panoramica del Progetto}

Il progetto esplora diversi approcci di apprendimento per rinforzo (RL) e ragionamento
basato su modelli linguistici (LLM/VLM) applicati all'ambiente
\texttt{MiniGrid-DoorKey} di Gymnasium. Il task richiede a un agente di raccogliere
una chiave, aprire una porta e raggiungere l'obiettivo verde, navigando in una
griglia bidimensionale con osservazione parziale o completa.

Il codice è organizzato in tre cartelle:
\begin{itemize}
    \item \texttt{agent/} — algoritmi di addestramento (Q-Learning, DDQN, LLM-guided);
    \item \texttt{env/} — wrapper e sistema di reward personalizzato;
    \item \texttt{prove/} — esperimenti e prototipi con modelli visivi (VLM/LLM).
\end{itemize}

%─────────────────────────────────────────────────────────────
\section{Moduli Ambiente (\texttt{env/})}

\subsection{reward-system (\texttt{env/rewardsystem.py})}

Implementa \texttt{DoorKeyRewardSystem}, un \texttt{gym.Wrapper} che sostituisce il
reward sparso originale con un segnale denso e strutturato.

Il reward totale per ogni step è la somma di cinque componenti configurabili
tramite la classe \texttt{RewardConfig}:

\begin{enumerate}
    \item \textbf{Reward ambientale} — valore originale restituito dall'ambiente.
    \item \textbf{Bonus milestone} — bonus una-tantum al raggiungimento di
          sotto-obiettivi chiave (raccolta chiave, apertura porta, goal).
    \item \textbf{Potential-based shaping} — $F(s,s') = \gamma\,\Phi(s') - \Phi(s)$
          dove $\Phi$ rappresenta il progresso normalizzato BFS verso l'obiettivo
          corrente lungo la sequenza \emph{chiave~→~porta~→~goal}.
    \item \textbf{Penalità regressione} — punisce il ritorno a fasi precedenti
          (es.\ perdere la chiave dopo averla raccolta).
    \item \textbf{Penalità temporale} — piccola costante negativa per step,
          per incentivare soluzioni rapide.
\end{enumerate}

Il progresso viene tracciato attraverso l'enum \texttt{Stage}
(\texttt{NO\_KEY} → \texttt{HAS\_KEY} → \texttt{DOOR\_OPEN} → \texttt{GOAL\_REACHED})
e un campo \texttt{curr\_progress} $\in [0,1]$ calcolato via BFS.

\subsection{vlm\_wrapper (\texttt{env/vlm\_wrapper.py})}

Implementa \texttt{VLMDebugWrapper}, un wrapper che arricchisce il reward
con la valutazione di un modello linguistico locale (Ollama, modello
\texttt{nemotron-3-nano:4b}).

Ogni \texttt{query\_every} step, il wrapper:
\begin{enumerate}
    \item Converte lo stato della griglia in formato CSV testuale
          (oggetti codificati come \texttt{wall}, \texttt{key\_yellow}, \texttt{door\_yellow\_locked}, ecc.);
    \item Affianca il frame passato e quello corrente;
    \item Invia un prompt strutturato al modello, che deve analizzare la scena,
          valutare l'ultima azione e restituire \texttt{reward} e \texttt{progress}
          in tag XML \texttt{<metrics>};
    \item Somma il reward VLM al reward ambientale.
\end{enumerate}

Il prompt include regole esplicite sull'ambiente, criteri di assegnazione del
reward e istruzioni di formato rigorose per facilitare il parsing della risposta.

%─────────────────────────────────────────────────────────────
\section{Agenti (\texttt{agent/})}

\subsection{doorkey-vanilla (\texttt{agent/doorkey\_vanilla.py})}

Baseline di Q-Learning tabellare sull'ambiente \texttt{MiniGrid-DoorKey-8x8}.
Lo stato è codificato come coppia di canali dell'osservazione immagine
(tipo oggetto e colore) serializzata in bytes. La politica $\epsilon$-greedy
usa un \textbf{cosine annealing} dell'esplorazione:
\[
  \epsilon(t) = \epsilon_{\min} + (\epsilon_{\max} - \epsilon_{\min})\,
  \frac{1 + \cos\!\left(\pi \left(\tfrac{t}{T}\right)^k\right)}{2}
\]
L'agente supporta tre modalità via CLI: \texttt{sweep} (ottimizzazione Bayesiana
degli iperparametri con WandB), \texttt{train} e \texttt{test}.
I valori iniziali ottimistici della Q-table ($Q_0 = 2$) favoriscono
l'esplorazione nelle fasi iniziali.

\subsection{doorkey-shaping (\texttt{agent/doorkey\_shaping.py})}

Estende il vanilla aggiungendo il \textbf{reward shaping potential-based} direttamente
nel wrapper \texttt{DoorKeyProgressReward}. La funzione potenziale $\Phi(s)$ è
calcolata tramite BFS sulla griglia, tenendo conto della sequenza di sotto-obiettivi
e della presenza della chiave. Il fattore di scala dello shaping è un iperparametro
incluso nello sweep Bayesiano, insieme a $\alpha$, $\gamma$, $\epsilon_{\min}$ e $k$.
Il training prevede 50\,000 episodi (contro i 25\,000 del vanilla), con il modello
salvato su disco per il test.

\subsection{doorkey-DDQN (\texttt{agent/doorkey-DDQN.py})}

Implementa un agente \textbf{Double DQN (DDQN) con Prioritized Experience Replay (PER)}.
La rete neurale (\texttt{DualHeadDDQN}) ha due rami in ingresso:
\begin{itemize}
    \item \textbf{Grid MLP} — processa i canali dell'osservazione immagine
          normalizzati (tipo oggetto e stato porta) tramite 4 strati lineari con
          \texttt{LayerNorm} e \texttt{ReLU};
    \item \textbf{Scalar MLP} — processa 9 feature scalari: direzione agente
          (one-hot 4 bit), fase corrente (one-hot 4 bit) e progresso continuo.
\end{itemize}
I due rami vengono concatenati e passati a un MLP finale che produce i Q-values.
Il target network viene aggiornato con soft update ($\tau = 0{,}001$) ogni step
o con hard copy ogni 4\,000 step. Il PER usa priorità proporzionali agli errori TD
per campionare le esperienze più informative. Il training è tracciato su WandB con
sweep Bayesiano degli iperparametri \texttt{lr}, \texttt{gamma}, \texttt{eps\_decay}.

\subsection{doorkey-llm2 (\texttt{agent/doorkey-llm2.py})}

Combina Q-Learning tabellare con il reward fornito da \texttt{VLMDebugWrapper}.
Lo stato è compresso in una tupla di cinque elementi
$(\Delta x,\, \Delta y,\, d,\, \text{progress\_bin},\, \text{stage\_idx})$
che codifica la distanza relativa all'obiettivo corrente, la direzione dell'agente,
il progresso discretizzato e la fase. Per i primi 150 episodi il VLM viene ignorato
(\texttt{ignore=True}) per permettere all'agente di esplorare liberamente; a seguire
il reward VLM viene integrato nel loop di training. L'agente supporta
valutazione finale e visualizzazione \texttt{human}.

%─────────────────────────────────────────────────────────────
\section{Prototipi e Esperimenti (\texttt{prove/})}

\subsection{provaLLM (\texttt{prove/provaLLM.py})}

Prototipo che pilota l'agente MiniGrid usando un LLM testuale remoto
(API Groq, modello multimodale). La griglia viene convertita in CSV testuale
e inviata al modello tramite due prompt alternativi (italiano e inglese).
Il modello restituisce reward, progresso e, in una variante, anche le coordinate
$(\Delta x, \Delta y)$ dell'obiettivo corrente, usando tag XML \texttt{<metrics>}
e \texttt{<analysis>}. Il file contiene anche la struttura per un agente Q-Learning
integrato con questo reward LLM.

\subsection{provaVLM6 (\texttt{prove/provaVLM6.py})}

Esperimento con LLM multimodale (Groq) che riceve una \textbf{sequenza di 5 frame
visivi} concatenati orizzontalmente come singola immagine. Il modello deve analizzare
la storia temporale per dedurre le regole dell'ambiente, identificare l'obiettivo
intermedio e restituire il reward e il livello di completamento.
Due prompt distinti gestiscono: (1) la selezione dell'azione successiva
con giustificazione testuale; (2) la valutazione della performance con metriche
numeriche.

\subsection{provaVLM2 (\texttt{prove/provaVLM2.py})}

Prototipo analogo che usa il modello \textbf{Qwen3-VL-4B-Instruct} in locale
tramite \texttt{transformers} e \texttt{torch.compile}. Riceve un singolo frame
visivo renderizzato dell'ambiente MiniGrid, costruisce il messaggio nel formato
atteso da Qwen-VL e chiede al modello di analizzare la scena e pianificare la
sequenza di azioni per risolvere il livello. Il \texttt{VLMDebugWrapper} interno
memorizza gli ultimi frame in una \texttt{deque} e interroga il modello ogni
\texttt{query\_every} step.

%─────────────────────────────────────────────────────────────
\section{Riepilogo}

\begin{table}[h]
\centering
\small
\begin{tabular}{llll}
\toprule
\textbf{File} & \textbf{Algoritmo} & \textbf{Reward} & \textbf{Modello} \\
\midrule
doorkey-vanilla    & Q-Learning tabellare & Sparso (env)       & — \\
doorkey-shaping    & Q-Learning tabellare & Potential shaping  & — \\
doorkey-DDQN       & DDQN + PER           & Sistema strutturato & Rete neurale \\
doorkey-llm2       & Q-Learning + VLM     & LLM (Ollama/locale) & nemotron-3-nano \\
reward-system      & Wrapper              & Denso multi-componente & — \\
vlm\_wrapper       & Wrapper              & LLM testuale       & nemotron-3-nano \\
provaLLM           & Prototipo Q+LLM      & LLM testuale remoto & Groq API \\
provaVLM6          & Prototipo VLM        & LLM visivo remoto  & Groq multimodale \\
provaVLM2          & Prototipo VLM        & VLM visivo locale  & Qwen3-VL-4B \\
\bottomrule
\end{tabular}
\caption{Panoramica dei file del progetto.}
\end{table}

Il progetto segue una progressione logica: si parte da una baseline Q-Learning
(\textit{vanilla}), si aggiunge reward shaping (\textit{shaping}), si introduce
una rete neurale profonda con memoria prioritizzata (\textit{DDQN}), e infine
si sperimenta la sostituzione o integrazione del reward con giudizi espressi da
modelli linguistici — prima via testo, poi via immagini — sia locali che remoti.

\end{document}
